% ============================================================================
% CAPÍTULO 1: INTRODUCCIÓN
% ============================================================================
% 
% INSTRUCCIONES GENERALES:
% - Este capítulo presenta una visión general de la investigación
% - Debe ser claro, conciso y motivar al lector
% - Extensión típica: 10-15 páginas
% - Reemplaza TODO el texto entre [corchetes]
% ============================================================================

\chapter{Introducción}
\label{cap:introduccion}

\justifying
\setlength{\parskip}{0pt}

% ----------------------------------------------------------------------------
% CONTENIDO REQUERIDO PARA INTRODUCCIÓN:
% ----------------------------------------------------------------------------
% Este capítulo debe incluir una visión general de:
% - El contexto y relevancia del tema de investigación
% - La problemática que abordas
% - Una breve descripción de lo que harás
% - La estructura del documento
%
% NOTA IMPORTANTE: La Delimitación del Objeto de Estudio (planteamiento del
%                  problema, pregunta de investigación, hipótesis, objetivos)
%                  va en el CAPÍTULO 3, NO aquí.
% ----------------------------------------------------------------------------

\noindent 
{[Escribe aquí tu introducción siguiendo esta estructura:]}

\textbf{Párrafo 1: Contexto General}

{[Introduce el tema general de tu investigación. Establece por qué es importante. Menciona datos epidemiológicos o estadísticas relevantes. Cita fuentes actuales.]}

\textbf{Párrafo 2: Problemática Específica}

{[Describe el problema específico que abordas. ¿Qué vacío en el conocimiento existe? ¿Qué necesidad práctica se debe resolver? Cita literatura que respalde la problemática.]}

\textbf{Párrafo 3: Aproximación o Solución Propuesta}

{[Describe brevemente cómo tu investigación aborda el problema. ¿Qué enfoque usas? ¿Qué hace único a tu trabajo? NO entres en detalles metodológicos (eso va en Cap. 5).]}

\textbf{Párrafo 4: Relevancia y Justificación}

{[Explica por qué tu investigación es relevante científica, metodológica y socialmente. ¿A quién beneficia? ¿Qué aporta al campo?]}

\textbf{Párrafo 5: Estructura del Documento (Opcional)}

{[Puedes incluir una breve descripción de cómo está organizada tu tesis capítulo por capítulo, o eliminarlo si lo prefieres.]}

% ----------------------------------------------------------------------------
% EJEMPLO DE REDACCIÓN:
% ----------------------------------------------------------------------------
%
% \noindent El comportamiento sedentario (CS) se ha consolidado como uno de los 
% principales retos de salud pública del siglo XXI, vinculado con...
% \cite{WHO2020,Bull2020}.
%
% Tradicionalmente, el estudio del CS se ha basado en instrumentos de 
% autoinforme como el GPAQ, que presentan limitaciones... \cite{Autor2023}.
%
% En este contexto, la presente investigación propone [tu solución/modelo]...
%
% La relevancia de este trabajo radica en... \cite{Referencia2024}.
%
% ----------------------------------------------------------------------------

% ============================================================================
% NOTAS FINALES:
% ============================================================================
% - Recuerda citar fuentes relevantes usando \cite{clave_referencia}
% - Mantén un tono académico pero accesible
% - Evita entrar en detalles metodológicos (eso va en Cap. 5)
% - No incluyas objetivos ni hipótesis aquí (eso va en Cap. 3)
% ============================================================================
